\SourceLeaf{32} oder, genauer gesagt, dass diese Summe niemals positiv werden kann. Sind also $\gamma,\gamma',\gamma''$ usw. von $c,c',c''$ usw. verschiedene, aber mit den Bedingungen des Systems vereinbare Orte und $\theta,\theta';\theta''$ usw. die Winkel, die $c\gamma,c'\gamma',c''\gamma''$ usw. mit $cb,c'b',c''b''$ usw. bilden, so ist $\Sigma m.cb.c\gamma.\cos\theta$ stets entweder $0$ oder negativ. Da nun
\[
\gamma b^2=cb^2+c\gamma^2-2cb.c\gamma.\cos\theta
\]
gilt, ist klar, dass
\[
\Sigma m.\gamma b^2-\Sigma m.cb^2=\Sigma m.c\gamma^2-2\Sigma m.cb.c\gamma.\cos\theta
\]
folglich immer positiv ist. Also ist $\Sigma m.\gamma b^2$ immer größer als $\Sigma m.cb^2$, das heißt, $\Sigma m.cb^2$ ist ein Minimum. W. Z. B. W.

Es ist sehr bemerkenswert, dass die Natur die freien Bewegungen, wenn sie mit den notwendigen Bedingungen nicht vereinbar sind, auf genau dieselbe Weise verändert, wie der rechnende Mathematiker nach der Methode der kleinsten Quadrate Beobachtungen ausgleicht, die sich auf Größen beziehen, die durch eine notwendige Abhängigkeit miteinander verbunden sind. Diese Analogie ließe sich noch weiter verfolgen, was jedoch nicht zu meiner gegenwärtigen Absicht gehört.
\par\medskip\centerline{\rule{.3\textwidth}{.3pt}}
